% Generated draft: requires mathematical and source review.
\documentclass[11pt]{article}
\usepackage[margin=25mm]{geometry}
\usepackage{fontspec}
\setmainfont{Latin Modern Roman}
\usepackage{amsmath,amssymb,mathtools,unicode-math}
\setmathfont{Latin Modern Math}
\usepackage{enumitem,needspace,xurl,hyperref}
\setlength{\parindent}{0pt}
\setlength{\parskip}{5pt}
\setlength{\emergencystretch}{3em}
\newcommand{\N}{\mathbb{N}}
\newcommand{\Z}{\mathbb{Z}}
\newcommand{\Q}{\mathbb{Q}}
\newcommand{\R}{\mathbb{R}}
\newcommand{\C}{\mathbb{C}}
\newcommand{\F}{\mathbb{F}}
\newcommand{\A}{\mathbb{A}}
\newcommand{\PP}{\mathbb{P}}
\newcommand{\E}{\mathbb{E}}
\newcommand{\eps}{\varepsilon}
\newcommand{\dd}{\,\mathrm{d}}
\begin{document}
\section*{UnsolvedMath ID: 700828}
\textbf{Source problem code:} catalogue entry
\subsection{Definitions and mathematical statement}
Let \(\mathcal V\) be a non-trivial variety of groups, meaning a class defined by group identities that contains a group with more than one element. Write \(F_m(\mathcal V)\) for the relatively free group of rank \(m\) in \(\mathcal V\). Fix an integer \(r>0\) and a free basis \(a_1,\ldots,a_r\) of \(F_r(\mathcal V)\).

The source does not define ordinary discrimination. Using its usual simultaneous homomorphism-separation meaning, say that \(F_r(\mathcal V)\) discriminates \(\mathcal V\) if, for every positive integer \(m\) and every finite nonempty set \(S\subseteq F_m(\mathcal V)\setminus\{1\}\), there exists a homomorphism \(h:F_m(\mathcal V)\to F_r(\mathcal V)\) such that
\[
 h(u)\ne1\qquad\text{for every }u\in S.
\]
This interpretation of the undefined source term requires review.

An inclusion of \(F_r(\mathcal V)\) as a varietally free factor means that its displayed basis extends to a free basis of the containing relatively free group. For a nonnegative integer \(t\), identify \(F_r(\mathcal V)\) with the subgroup generated by \(a_1,\ldots,a_r\) in the relatively free group
\[
 F_{r+t}(\mathcal V)=F_{\mathcal V}(a_1,\ldots,a_r,b_1,\ldots,b_t).
\]
Following the source, \(F_r(\mathcal V)\) strongly discriminates \(\mathcal V\) if the following holds for every integer \(k\ge0\), every integer \(n\ge1\), and every family of group words \(w_1,\ldots,w_n\) in the displayed parameters and variables: whenever, for some \(t\ge0\), there exists \(\mathbf c\in F_{r+t}(\mathcal V)^k\) satisfying
\[
 w_i(a_1,\ldots,a_r,\mathbf c)\ne1\qquad(1\le i\le n),
\]
there exists \(\mathbf d\in F_r(\mathcal V)^k\) satisfying
\[
 w_i(a_1,\ldots,a_r,\mathbf d)\ne1\qquad(1\le i\le n).
\]
The parameters \(a_j\) remain fixed throughout. Although the source does not specify whether the containing free group may have infinite rank, finite ranks suffice here: a finite solution tuple involves only finitely many additional basis elements.

The group-theoretic question is whether there exist a non-trivial variety \(\mathcal V\) and an integer \(r>0\) such that
\[
 F_r(\mathcal V)\text{ discriminates }\mathcal V
 \quad\text{but does not strongly discriminate }\mathcal V.
\]
Thus failure of strong discrimination must be witnessed by a finite system of inequalities with these fixed free parameters that is solvable in a relatively free extension but has no solution in \(F_r(\mathcal V)\).

The source also asks the analogous existence question for general algebras in universal algebra. A usual interpretation replaces nonidentity group elements by pairs of distinct elements of relatively free algebras, discrimination by simultaneous separation of finitely many such pairs by a homomorphism, and the strong condition by transfer of finite systems of term inequalities with fixed free parameters from a free extension to the original free algebra. The source supplies neither a signature nor an explicit definition of this algebraic analogue; this description records a candidate interpretation, and the exact intended scope remains unresolved.
\subsection{Short English statement}
Does some non-trivial group variety have a relatively free group of positive finite rank that discriminates the variety, yet fails to realize a finite system of inequalities with fixed free-generator parameters that is realizable in a relatively free extension? The source also asks for an analogue for general algebras, without specifying its precise scope. Ordinary discrimination is interpreted in its usual simultaneous separation sense.
\subsection{Sources}
\begin{itemize}
\item Supplied local UnsolvedMath registry record, permanent ID \(700828\), exact source code \texttt{catalogue entry}, title \emph{Kourovka Notebook Problem 12.8}. The supplied status is \texttt{open}; this is a registry label, not an independently verified mathematical or literature assessment.
\item \emph{The Kourovka Notebook (Unsolved Problems in Group Theory)}, Problem 12.8, 12th issue (1992); the supplied background also identifies the 21st issue edition (2026), \texttt{arXiv:1401.0300v48}. Proposed by A. M. Gaglione and D. Spellman in 1992. The supplied collection description names E. I. Khukhro and V. D. Mazurov as editors.
\item The Notebook attribution and version identifier are supplied evidence. No external source was visited or independently checked, and no source-verification date or external link was supplied.
\end{itemize}
\end{document}
